\documentclass[10pt]{article}
\usepackage[T1]{fontenc}
\usepackage[utf8]{inputenc}
\usepackage[french]{babel}
\usepackage{amsmath,amssymb,mathrsfs,mathpazo}
\usepackage{geometry}
\geometry{a4paper,margin=2cm}
\usepackage{graphicx,array,tabularx,tabularray,colortbl,multirow,multicol,booktabs,enumitem,pifont,ulem,fancybox,diagbox,cancel}
\usepackage{pstricks,pstricks-add,pst-plot,pst-tree,pst-node,pst-text,pst-func,pst-math,pst-bspline,pst-3dplot,pst-eucl,pst-solides3d}
\let\pscalculate\undefined
\usepackage{pst-bezier}
\usepackage{tikz,pgfplots,tkz-tab}
\usetikzlibrary{arrows,arrows.meta,patterns,calc,positioning,trees,decorations.markings,shapes.geometric,shapes.misc,shapes.arrows,intersections,decorations.pathmorphing,decorations.pathreplacing,backgrounds,shadows,babel}
\pgfplotsset{compat=1.18}
\usepgfplotslibrary{fillbetween}
\usepackage[np]{numprint}
\usepackage{listings,makecell,pas-tableur,twoopt}
\let\barre\undefined
\usepackage[pstricks]{tablvar}
\definecolor{darkgreen}{RGB}{0,100,0}
\providecommand{\red}{\color{red}}
\providecommand{\blue}{\color{blue}}
\providecommand{\green}{\color{green}}
\providecommand{\black}{\color{black}}
\providecommand{\white}{\color{white}}
\providecommand{\gray}{\color{gray}}
\providecommand{\cyan}{\color{cyan}}
\providecommand{\magenta}{\color{magenta}}
\providecommand{\yellow}{\color{yellow}}
\providecommand{\violet}{\color{violet}}
\providecommand{\orange}{\color{orange}}
\providecommand{\euro}{\texteuro}
\providecommand{\EUR}{\texteuro}
\providecommand{\ieme}{\textsuperscript{e}}
\providecommand{\degres}{\ensuremath{^\circ}}
\providecommand{\degree}{\ensuremath{^\circ}}
\providecommand{\vv}[1]{\overrightarrow{#1}}
\providecommand{\Coord}[1]{#1}
\providecommand{\decosix}{}
\providecommand{\starredbullet}{\ensuremath{\bullet}}
\pagestyle{empty}
\setlength{\parindent}{0pt}
\usepackage{adjustbox}\DeclareUnicodeCharacter{2236}{:}\DeclareUnicodeCharacter{2190}{\ensuremath{\leftarrow}}
\providecommand{\R}{\mathbb{R}}
\providecommand{\N}{\mathbb{N}}
\providecommand{\Z}{\mathbb{Z}}
\providecommand{\Q}{\mathbb{Q}}
\providecommand{\D}{\mathbb{D}}
\providecommand{\C}{\mathbb{C}}
\providecommand{\e}{{\mathrm{e}}}
\providecommand{\vect}[1]{\overrightarrow{#1}}
\providecommand{\vectt}[1]{\overrightarrow{\mathrm{#1}}}
\providecommand{\vvt}[1]{\overrightarrow{\mathrm{#1}}}
\providecommand{\np}[2][]{{#2~#1}}
\providecommand{\vv}[1]{{\overrightarrow{#1}}}
\renewcommand{\d}{{\mathrm{d}}}
\providecommand{\degres}{{^\circ}}
\providecommand{\degre}{{^\circ}}
\providecommand{\blue}{\color{blue}}
\providecommand{\red}{\color{red}}
\providecommand{\enskip}{\quad}
\providecommand{\llbracket}{\lbrack\!\lbrack}
\providecommand{\rrbracket}{\rbrack\!\rbrack}
\providecommand{\euro}{{\text{€}}}
\providecommand{\no}{{\text{n°}}}
\providecommand{\interval}[2]{{[#1;#2]}}
\providecommand{\intervFO}[2]{{[#1;#2[}}
\providecommand{\dotfill}{\ldots}

\providecommand{\ds}{\displaystyle}
\providecommand{\textding}[1]{#1}
\providecommand{\Oij}{\(\mathrm{O};\vec{\imath},\vec{\jmath}\)}
\providecommand{\Oijk}{\(\mathrm{O};\vec{\imath},\vec{\jmath},\vec{k}\)}
\providecommand{\Ouv}{\(\mathrm{O};\vec{u},\vec{v}\)}
\providecommand{\intervFF}[2]{\ensuremath{[#1;#2]}}
\renewcommand{\euro}{\ifmmode\text{\texteuro}\else\texteuro\fi}
\DeclareUnicodeCharacter{00B0}{\ifmmode^{\circ}\else\textdegree\fi}
\DeclareUnicodeCharacter{20AC}{\ifmmode\text{\texteuro}\else\texteuro\fi}
% Énoncés : Éducation nationale. Transcription et corrigés : APMEP, Denis Vergès et contributeurs. Découpage et classement : CoopMaths, Sébastien Lozano, Éric Elter et contributeurs. Adaptation de la présentation : Maths Explorer.
% CC BY-SA : https://coopmaths.fr/www/about/
% https://coopmaths.fr/alea/static/bac/2024/tex/bac_2024_09_sujet1_polynesie_3_cor.tex
\begin{document}
\noindent\textit{APMEP / CoopMaths. CC BY-SA. Présentation adaptée par Maths Explorer.}\par\medskip
\medskip

On considère une pyramide à base carrée formée de boules identiques empilées les unes sur les autres:

\begin{list}{\textbullet}{}
\item le 1\up{er} étage, situé au niveau le plus haut, est composé de 1 boule ;
\item le 2\up{e} étage, niveau juste en dessous, est composé de 4 boules ;
\item le 3\up{e} étage possède 9 boules ;
\item \ldots
\item le $n$-ième étage possède $n^2$ boules.
\end{list}

\medskip

Pour tout entier $n \geqslant 1$, on note $u_n$ le nombre de boules qui composent le $n$-ième étage en partant du haut de la pyramide. Ainsi, $u_n = n^2$.

\begin{enumerate}
\item Le nombre total de boules d'une pyramide de 4 étages est:\\
$u_1+u_2+u_3+u_4=1^2+2^2+3^2+4^2=1+4+9+16=30$.

\item On considère la suite $\left(S_n\right)$ définie pour tout entier $n \geqslant 1$ par 
$S_n = u_1 + u_2 + \ldots + u_n.$

	\begin{enumerate}
		\item $S_5 = u_1+u_2+u_3+u_4+u_5  = \left ( u_1+u_2+u_3+u_4 \strut\right ) +u_5 30+5^2 = 55$
		
Le nombre total de boules d'une pyramide de 5 étages est: $S_5=55$.
		 		 
		\item
		
On complète la fonction \texttt{pyramide} ci-dessous de sorte que, pour tout entier naturel non nul $n$, l'instruction \texttt{pyramide(n)} renvoie le nombre de boules composant une pyramide de $n$ étages.

\begin{center}
\ttfamily
\begin{tabular}{|l|}\hline
def pyramide(n) :\\
\qquad S = 0\\
\qquad for i in range(1,~n+1) :\\
\qquad \qquad S =  \blue S + i**2\\
\qquad return \blue S\\ \hline
\end{tabular}
\end{center}

		\item 
\begin{list}{\textbullet}{Pour tout entier naturel $n$ :}
\item 
$\begin{aligned}[t]
\dfrac{n(n+1)(2n+1)}{6} + (n+1)^2 
& = \dfrac{(n+1)(2n^2+n)}{6} + \dfrac{6(n+1)^2}{6}\\
& = \dfrac{(n+1)(2n^2+n+6n+6)}{8}\\
 & = \dfrac{(n+1)(2n^2+7n+6)}{6}
 \end{aligned}$
\item 
$\begin{aligned}[t]
\dfrac{(n + 1)(n +2)[2(n + 1) + 1]}{6}
& = \dfrac{(n+1)(n+2)(2n+3)}{6}
= \dfrac{(n+1)(2n^2+3n+4n+6)}{6}\\
& = \dfrac{(n+1)(2n^2+7n+6)}{6}
 \end{aligned}$
\end{list}		
		
Donc $\dfrac{n(n+1)(2n+1)}{6} + (n+1)^2 = \dfrac{(n + 1)(n +2)[2(n + 1) + 1]}{6}$

		\item On démontre par récurrence que pour tout entier $n \geqslant 1$ : 
$S_n = \dfrac{n(n + 1)(2n + 1)}{6}$.

\begin{list}{\textbullet}{}
\item \textbf{Initialisation}

Pour $n=1$, on a $S_1=1$ et $\dfrac{n(n + 1)(2n + 1)}{6} = \dfrac{1(1 + 1)(2\times 1 + 1)}{6} = \dfrac{1\times 2\times 3}{6} =1$

Donc la propriété est vraie au rang 1.

\item \textbf{Hérédité}

On suppose la propriété vraie au rang $n$, avec $n \geqslant 1$, c'est-à-dire $S_n = \dfrac{n(n + 1)(2n + 1)}{6}$. C'est l'hypothèse de récurrence.

$\begin{aligned}[t]
S_{\blue n+1} 
& = u_1+u_2+\cdots +u_n+u_{n+1} =  \left (u_1+u_2+\cdots +u_n \strut\right )+u_{n+1} \\
& = S_n+(n+1)^2 =  \dfrac{n(n + 1)(2n + 1)}{6}+(n+1)^2 \\
& = \dfrac{(n + 1)(n +2)[2(n + 1) + 1]}{6} \text{ (d'après la question précédente)}\\
& = \dfrac{({\blue n + 1})[({\blue n + 1}) + 1][2({\blue n + 1}) + 1]}{6}
\end{aligned}$

Donc la propriété est vraie au rang ${\blue n + 1}$.

\item \textbf{Conclusion}

La propriété est vraie au rang 1, et elle est héréditaire pour tout $n\geqslant 1$. \\
Donc, d'après le principe de récurrence, la propriété est vraie pour tout $n\geqslant 1$.
\end{list}

On a donc démontré par récurrence que pour tout entier $n \geqslant 1$ : 
$S_n = \dfrac{n(n + 1)(2n + 1)}{6}$.
	\end{enumerate}
	
\item Un marchand souhaite disposer des oranges en pyramide à base carrée. Il possède $200$~oranges.

Il faut donc trouver le plus grand entier $n$ tel que $S_n \leqslant 200$.

On calcule:
$S(7)=\dfrac{7(7+1)(2\times 7+1}{6}=140<200$ et
$S(8)=\dfrac{8(8+1)(2\times 8+1)}{6}=204>200$.

Le marchand utilise donc 140 oranges pour construire une pyramide à 7 étages.
\end{enumerate}

\bigskip

\end{document}
